% Source: Wald GR, physical PDF p. 24
%         (printed p. 11).
% Coverage: Chapter 2 opening material before Section 2.1.
% Figures/tables/footnotes: none.
% Status: source-reviewed and compile-clean.
% Uncertainties: none.

In this chapter we lay the foundations for a precise, mathematical formulation of
general relativity by obtaining some basic properties of manifolds and tensor fields.
As defined in Section~\ref{sec:2.1}, an \(n\)-dimensional manifold is a set that has
the local differential structure of \(\mathbb{R}^n\) but not necessarily its global
properties.  In Section~\ref{sec:2.2} we define tangent vectors as directional
derivative operators acting on functions defined on a manifold.  We obtain there
some important properties of coordinate bases of the tangent space and tangents to
curves.  Tensors are introduced in Section~\ref{sec:2.3}, and the notion of a
metric is defined.  Finally, in Section~\ref{sec:2.4} we introduce the abstract
index notation for tensors, which we shall use throughout the remainder of this
book.  We will use a fair number of standard mathematical symbols in this chapter,
and the reader unfamiliar with these symbols should consult the section ``Notation
and Conventions'' at the beginning of this book.
